%-------------------------
% Resume in Latex
% Author : Sourabh Bajaj
% Modified by: Matthew Knight
% License : MIT
%------------------------

\documentclass[letterpaper,10pt]{article}

\usepackage[empty]{fullpage}
\usepackage{titlesec}
\usepackage[usenames,dvipsnames]{color}
\usepackage{enumitem}
\usepackage[hidelinks]{hyperref}
\usepackage{fancyhdr}
\usepackage{tabularx}
\usepackage{graphicx}

\graphicspath{{images/}}

\pagestyle{fancy}
\fancyhf{} % clear al/l header and footer fields
\renewcommand{\headrulewidth}{0pt}
\renewcommand{\footrulewidth}{0pt}

% Adjust margins
\addtolength{\oddsidemargin}{-0.5in}
\addtolength{\evensidemargin}{-0.5in}
\addtolength{\textwidth}{1in}
\addtolength{\topmargin}{-.5in}
\addtolength{\textheight}{1.0in}

\raggedbottom
\raggedright
\setlength{\tabcolsep}{0in}
\setlength{\footskip}{5pt}

% Sections formatting
\titleformat{\section}{
	\vspace{-4pt}\scshape\raggedright\large
}{}{0em}{}[\color{black}\titlerule \vspace{-5pt}]

%-------------------------
% Custom commands
\newcommand{\resumeItem}[2]{
    \item\small{
    \textbf{#1}{: #2 \vspace{-2pt}}
  }
}

\newcommand{\ResumeSubheading}[4]{
  \vspace{-1pt}\item
    \begin{tabular*}{\textwidth}[t]{l@{\extracolsep{\fill}}r}
      \textbf{#1} & #2 \\
      \textit{\small#3} & \textit{\small #4} \\
    \end{tabular*}
  \vspace{-5pt}
}

\newcommand{\resumeSubItem}[2]{\resumeItem{#1}{#2}\vspace{-4pt}}

\newenvironment{SubheadingList}
  {\begin{itemize}[label=,leftmargin=0em]}
  {\end{itemize}}

\newenvironment{ItemList}
  {\begin{itemize}[label=,leftmargin=1em]}
  {\end{itemize}\vspace{-5pt}}

%-------------------------------------------
%%%%%%  CV STARTS HERE  %%%%%%%%%%%%%%%%%%%%%%%%%%%%


\begin{document}

%----------HEADING-----------------
\begin{tabularx}{\textwidth}{X c >{\raggedleft\arraybackslash}X}
  \begin{tabular}{l}
    1766 Biatecki Rd. \\
    Revelstoke, BC \\
    V0E 2S1, Canada
  \end{tabular} &
  \textbf{\href{http://mknight.io}{\LARGE Matthew Knight}} &
  \begin{tabular}{r}
    +1-778-791-4069 \\
    \href{mailto:matt@mknight.io}{matt@mknight.io}
  \end{tabular} \\ &
  \href{http://mknight.io}{\includegraphics[width=1.5em]{mknight.pdf}}
   \href{https://github.com/matt1795}{\includegraphics[width=1.5em]{github}}
  \href{https://gitlab.com/matt1795}{\includegraphics[width=1.5em]{gitlab}}
  \href{https://www.linkedin.com/in/matthew-knight-a9b70894/}{\includegraphics[width=1.5em]{linkedin}}
\end{tabularx}
%
\vspace{1em}

% About me:
An Embedded Software Developer interested in Industrial Automation, Distributed
Systems, and Real-Time Applications. I'm looking for a remote contract or
full-time position in order to live in Revelstoke, BC.

%--------PROGRAMMING SKILLS------------
\section{Skills \& Tools}
\begin{SubheadingList}
  \resumeItem{Proficient}{C, C++, Python, Git, Unix, Bash, Assembly, JavaScript,
  Verilog}
  \resumeItem{Familiar}{Go, SQL, Java}
\end{SubheadingList}

%-----------EXPERIENCE-----------------
\section{Experience}
\begin{SubheadingList}
  \ResumeSubheading
    {Software Developer II}{Oct 2018 -- Aug 2019}
    {Avigilon}{Vancouver, BC}
  \begin{ItemList}
    \resumeItem{C++ Visual Recognition Library}
      {Implemented debugging system to dump inputs/outputs, latency and
      throughput of pipeline that processed images using GPU hardware which
      verified the performance requirements of video analytics library.}
    \resumeItem{Convolutional Neural Network (CNN) Optimization}
      {Wrote utilities in C++ to take trained neural net models in Caffe or
      TensorFlow format and optimize them to a graphics processing unit (GPU)
      architecture which increased throughput 100-200\%.}
    \resumeItem{Certificate Management}
      {Used Yocto Linux to add certificate management to Avigilon's internal
      embedded Linux distribution which filled requirements of high security
      applications.}
    \resumeItem{Leveraged Knowledge}{in C, C++, Application Binary Interfaces,
    GCC, GDB, OpenCV, Software Debugging, Conan, CMake, Git, Template
    Metaprogramming, Docker, Linux, Protobuf, CNNs, Python, TensorRT,
    TensorFlow, Continuous Integration (CI), GPUs, Yocto, SSL Certificates, HTTP
    APIs, Go, Networking, React, Build Systems, Bash}
  \end{ItemList}

  \ResumeSubheading
    {Full-Stack Engineer}{Oct 2017 -- Oct 2018}
    {Plasmatic Technologies Inc.}{Vancouver, BC}
  \begin{ItemList}
    \resumeItem{Firmware}
      {Reverse engineered a proprietary serial communication protocol
      for a home security system by programming a microcontroller (MCU) with C to
      handle hardware interface and streamed data over a UART to a PC
      which analyzed packet structure using Python.}
    \resumeItem{Embedded Linux}
      {Built the linux kernel with custom drivers for a touch screen display
      Implemented a demo graphical user interface (GUI) using a Python graphics
      library called pygame.}
    \resumeItem{Back-End}
      {Became adept at using Node.js and some AWS modules (Cognito, DynamoDB,
      Iot-core, Lambda) in order to develop back-end for Plasmatic.}
    \resumeItem{Web and App Front-End}
      {Worked with Graphical Designer to implement part of the Web-based UI, and
      large components of the Mobile App using React and React Native
      respectively.}
    \resumeItem{Leveraged Knowledge}
    {in Microcontrollers, C, USART, Python, Serial Protocols, Reverse
    Engineering, Oscilloscope, Linux Kernel, Device Drivers, Node.js,
    JavaScript, AWS, NoSQL, MQTT, GraphQL, APIs, Mobile App Development,
    React-Native, Java, User Interfaces, User Experience, Git}
  \end{ItemList}

  \ResumeSubheading
    {R\&D Technician (8 Month Co-op and Summer Position)}{Jan 2015 -- Aug 2017}
    {Argus Controls Ltd.}{Surrey, BC}
  \begin{ItemList}
    \resumeItem{Retrofitting a Growth Chamber}
      {Replaced the control system of a germination chamber with Argus's, tuned
      the temperature and humidity controllers, and added necessary wiring to
      place testing sensors in chamber.}
    \resumeItem{QA}
      {Validation and verification of sensor accuracy performance using
      retrofitted growth chamber.}
    \resumeItem{Scripting}
      {Used Python to automate the data analysis of sensor performance.}
    \resumeItem{Leveraged Knowledge}
    {in Circuit Analysis, Control Systems, Validation and Verification, Data
    Analysis, Python, Statistics, Technical Documentation, and Thermodynamics}
  \end{ItemList}
\end{SubheadingList}

%-----------EDUCATION-----------------
\section{Education}
\begin{SubheadingList}
  \ResumeSubheading
    {Bachelor in Electrical Engineering}{Jan 2014 -- Jun 2018}
    {British Columbia Institute of Technology}{Burnaby, BC}
  \begin{ItemList}
    \resumeItem{Coursework}
    {Real-Time Embedded Systems, Digital Image and Video Processing, Wireless
    System Design, Software Systems, Feedback Control, Computer Networking,
    Digital System Design, Signal Processing and Filters, Sensors for
    Measurement and Control, Microcontroller Systems.}
  \end{ItemList}
\end{SubheadingList}

%-----------PROJECTS-----------------
\section{Projects}
\begin{SubheadingList}
  \ResumeSubheading
    {Prosthetics and Orthopedics Monitoring System (POMS)}{Sept 2017 -- Jun 2018}
    {Team Member}{Capstone}
    \begin{itemize}
      \item Collaborated with orthopedics clinic to determine requirements of
        Scoliosis brace monitoring system.
      \item Designed a communication protocol utilizing Bluetooth, End-to-end
        encryption in order to minimize power consumption of wearable sensor.
      \item Designed software backend to communicate with patient devices and
        provide clinicians with graphical interface for analyzing wear-time.
      \item Built prototype sensor that utilized capacitive sensing to measure
        wear-time and software backend for demonstration.
    \end{itemize}
  \ResumeSubheading
    {BCIT Sailbot}{Sept 2017 -- Jun 2018}
    {Organizer/Lead}{Extracurricular}
    \begin{itemize}
      \item Initiated club at BCIT to build and design an autonomous sailing
        robot for competition.
      \item Broke up project and collaborated with other students to research
        and design solutions for different aspects of the project.
      \item Designed software architecture for autonomous navigation.
      \item After two years, prototypes for a fixed-wing sail, hull, and novel
        wind sensor built.
    \end{itemize}
  \ResumeSubheading
    {multiformats}{Sept 2019}
    {Maintainer}{Personal}
    \begin{itemize}
      \item C++17 implementation of a collection of self-describing protocols
        for hashing, base encoding, network addressing, protocol negotiation.
      \item Analyzed requirements and created unit tests which prove
        interoperability with other library implementations
    \end{itemize}
  \ResumeSubheading
    {thermistor}{July 2019 -- Aug 2019}
    {Maintainer}{Personal}
    \begin{itemize}
      \item Utilized Template Metaprogramming to create a header-only library
        that generates lookup-tables for thermistors.
      \item Implemented unit tests that proved accuracy of mathematical models.
      \item Created an interface that makes it easy for users to work with
        different datasheet specifications and circuit configurations.
    \end{itemize}
  \ResumeSubheading
    {register-alias}{July 2019}
    {Maintainer}{Personal}
    \begin{itemize}
      \item Created method for safe access of special-function registers of ARM
        Cortex-M MCUs.
      \item Wrote python script that took standardized, XML-based register maps
        of MCUs and generated C++ code to represent each MCU as a type.
      \item Tested new method on hardware and prototyped Template
        Metaprogramming based peripheral drivers which safely and efficiently
        handle different register maps.
    \end{itemize}
  \ResumeSubheading
    {cold\_war\_chat}{Sept 2017}
    {Partner}{Curricular Project}
    \begin{itemize}
      \item Designed a communication stack for a real-time operating System (RTOS)
        which included AES-128 encryption and a cheap, half-duplex,
        on-off-keying (OOK) transciever.
      \item Configured each layer in the communication stack as a task within
        TI-RTOS
    \end{itemize}
  \ResumeSubheading
    {composit-video-driver}{Mar 2017}
    {Partner}{Curricular Project}
    \begin{itemize}
      \item Designed hardware system to take the frame buffer of a display
        output an analog composite video signal to drive an old cathode-ray-tube
        monitor.
      \item Utilized TDD to implement hardware design in Verilog
      \item Built a prototype board to interface with an FPGA development board
        and contained the Digital-to-Analog circuitry needed to generate the
        composite video signal.
    \end{itemize}
\end{SubheadingList}
%-------------------------------------------
\end{document}
